\subsection{Interpretation of Findings}

Our results demonstrate that computational routing structure—induced via attention sparsity manipulation—gives rise to observable patterns in population-level signals after dimensionality reduction. This supports the hypothesis that EEG reflects projections of internal pathway characteristics.

Three findings are notable:

\subsubsection{1. Strong Pathway-Sparsity Link ($r = 0.68$)}

Pathway metrics (routing entropy, path competition, inter-head divergence, etc.) form a coherent signature of attention structure. This validates our metric extraction and suggests that attention patterns are not random but carry task-relevant information.

\subsubsection{2. Partial Information Preservation ($r_{\mathcal{E}} = 0.42$)}

Linear projection to 105 channels with Gaussian noise reduces correlation from 0.68 to 0.42 (38\% retention rate). This is non-trivial: many mechanisms could completely destroy the signal. The fact that 18\% of variance is still explained suggests that routing structure maps to low-dimensional subspaces that survive noise.

\subsubsection{3. Meaningful Transfer ($\text{TSR} = 0.62$)}

Models trained on pathway features achieve 62\% of pathway-level performance when applied to EEG. This is substantially better than random (TSR $\approx 0$) and indicates a genuine causal relationship: pathways determine EEG patterns, not vice versa.

\subsection{Implications for EEG Interpretation}

\subsubsection{Neural Coding Hypothesis}

Our findings support a \textit{pathway coding} view of EEG: rather than directly encoding stimulus or task variables, EEG reflects the structure of information routing in neural circuits. This is consistent with theories emphasizing the importance of attentional routing \cite{Desimone1990, Luck2004} and efficient coding \cite{Barlow1961}.

\subsubsection{From Neural Networks to Neuroscience}

Artificial neural network models are increasingly used to predict brain activity \cite{Schrimpf2018, Conwell2022}. Our work suggests that mechanistic interpretability of these models—understanding how information is routed—may provide insights into neural coding principles observable in EEG.

\subsubsection{Inverse Problem: Inferring Routing from EEG}

Conversely, if EEG reflects routing patterns, then EEG-based decoding of cognitive states may be interpreted as decoding the underlying routing structure. This could improve BCI applications and clinical monitoring.

\subsection{Comparison to Prior Work}

Recent work on attention in transformers has shown that:

\begin{enumerate}
    \item Attention heads specialize into different roles (e.g., positional, syntactic) \cite{Clark2019}
    \item Attention patterns predict behavioral errors in humans \cite{Koyamada2019}
    \item Sparse attention reduces computational cost while preserving performance \cite{Child2019}
\end{enumerate}

Our contribution is to link attention patterns to population-level signals, demonstrating that routing structure survives a realistic measurement process (linear spatial mixing + noise).

\subsection{Limitations}

See Section~\ref{sec:limitations} for discussion of limitations.
